\chapter[Appendix: Data-driven estimation of symbolic representations of egocentric vision for the assessment of dysexecutive syndromes]{Data-driven Estimation of \\Symbolic Representations of Egocentric Vision \\ for the Assessment of Dysexecutive Syndromes \\ Appendix}
\label{appendix:part_2}



\section{Dataset preparation and preprocessing}
\label{apd:chapter_4_dataset_preparation}

\subsection{Recruitment}
\label{apd:chapter_4_methods}

% Included patients are recruited from dysexecutive patients already under medical care at the neurology department of the HIA Percy. %For this reason, we  anticipate that most will likely suffer from either TBI or Radiation-Induced Leukoencephalopathy (RIL), which are by far the most common conditions associated with DS in this department. 
% Exclusion criteria included motor impairments, severe visual or auditory deficits, and the presence of pre-existing neurodegenerative or psychiatric disorders affecting executive functions.

Included patients were recruited among individuals already under medical care at the neurology department of the HIA Percy. Specifically, brain-damaged participants had to present a dysexecutive syndrome diagnosed and monitor by the department’s neurology team. Eligibility criteria were defined to ensure the safety of participants and the clinical relevance of collected data. All participants, regardless of group, were required to be affiliated with a French social security system, aged over 18, and capable of providing informed consent. 

Exclusion criteria included pregnancy or breastfeeding, inability to read or understand French, and being under legal protection (“sauvegarde de justice”\footnotemark[1]).
\footnotetext[1]{”Sauvegarde de justice” refers to a French type of judicial protection of incapable adults.} Control participants were additionally excluded if they had any progressive or unstable neurological or psychiatric conditions. For brain-damaged patients, further exclusion criteria included any progressive psychiatric disease, except for post-traumatic stress disorder.

% Target sample sizes for the study were estimated to be sufficient for robust esstimation of the sensitivity and specificity of of our methodology to differentate patients with dyseecutive syndromes from the control group, at level to meet XX, respectlively  \cite{Obuchowski1998}


%CODE: include
% \subsection{Clinical data collection}
% \label{apd:chapter_4_clinical_measures}


% The aim of this study is to develop a digital phenotyping approach to assess sensorimotor and executive dysfunctions across the following five dimensions: static balance, dynamic balance, and locomotion; neurovegetative and emotional state; fine motor skills; visual information processing; and planning, initiation, mental flexibility, and memory. These syndromes can result from various neurological conditions, including traumatic brain injury (TBI) and lesions in the prefrontal cortex (LPR). Finally, patients’ EF will also be assessed using standard neuropsychological tests.  These tests will vary depending on what is appropriate for the usual care of each  patient, and thus will not necessary be the same for each patient, nor for every visit of the same patient. However, they should always include the Behavioral Dysexecutive  Syndrome Inventory (BDSI) [11], and the Dysexecutive Questionnaire (DEX) [12]  at Y1 and Y3


% \cite{Freitas2013} MOCA

% \cite{Godefroy2010} Dysexecutive syndrome: Diagnostic criteria and validation study,


% Dysexecutive behavioral changes were assessed using the Behavioral Dysexecutive Syndrome Inventory (BDSI), a structured interview of an informant assessing changes relative to previous behavior in twelve domains: (1) hypoactivity with apathy-abulia; (2) difficulties in anticipation, planning and initiation of activities; (3) disinterest and indifference to his/her own concern and others; (4) hyperactivity-distractibility-psychomotor instability; (5) irritability-impulsivity-aggressiveness; (6) euphoria, emotional lability and moria; (7) stereotyped and perseverative behavior; (8) environmental dependency; (9) anosognosia-anosodiaphoria; (10) spontaneous confabulations; (11) social behavior disorders; and (12) disorders of sexual, eating and urinary behavior [24].

% The diagnoses made by a multiprofessional conciliation meeting of neurologists working at the HIA Percy, leading to the diagnosis of DS and the inclusion or exclusion of patients in this study, will serve as the gold standard for diagnosing DS.  

% \cite{Roussel2016} identify IDDC as good 


\subsection{Classification of errors used in the manual annotations}
\label{apd:chapter_4_error_type_description}

We present a classification and examples of (i) descriptive errors and (ii) neuropsychological errors made by patients with dysexecutive syndromes in \Cref{fig:chapter_4_error_type_description_1} and \Cref{fig:chapter_4_error_type_description_2}, respectively \cite{Chevignard2008, Poncet2014}. These classifications of errors aim to provide relevant information to clinicians for specific rehabilitation intervention, and are typically used to evaluate executive functioning during the cooking task \cite{Poncet2014}.

% Poncet: Discriminant validity of the task was excellent, as
% patients made significantly more errors (mean total number of errors¼
% 107.5; SD¼ 74.3 versus 18.3; SD¼ 7.8)

% de Reliability of the Cooking Task in
% adults with acquired brain injury
% The Cooking Task is able to identify (even
% mild) EF disorders in patients with ABI (Chevignard et al., 2000). The
% error analysis used in the Cooking Task was developed following Schwartz
% et al.’s model (1998) of error analysis of task performance that considers
% different types of errors (“Additions”, “Omissions”, etc.), providing pertinent
% information to clinicians about the choice of rehabilitation intervention.


% Other error: taking the cake out of the oven with unprotected hands (poncet).

\begin{figure}[!htbp]
    \centering
    \includegraphics[width=1\linewidth]{chapters/chapter 5/figures/chapter_4_error_type_description_1.png}
    \caption{\textbf{Classification of errors used in the manual annotations.}}
    \label{fig:chapter_4_error_type_description_1}
 \end{figure}


\begin{figure}[!htbp]
    \centering
    \includegraphics[width=1\linewidth]{chapters/chapter 5/figures/chapter_4_error_type_description_2.png}
    \caption{\textbf{Classification of errors used in the manual annotations.}}
    \label{fig:chapter_4_error_type_description_2}
 \end{figure}

While these can be used to guide the analysis of participant's behavior during the cooking task, they define challenging targets that require further work to be automatically inferred from the symbolic representations of behavior. 

% \subsection{Dataset quality assessment and cleaning (supplementary material)}


% \subsection{Exhaustive video-level quality inspection}

% L'état actuel du dataset comprend:

% - Des videos à l'envers
% - Des inversions de conventions de placements des GoPros
% - Une amplitude de positionnement des GoPros qui varie beaucoup.

% Je propose de remédier à ces problèmes par inspection visuelle de l'ensemble du dataset, en utilisant la frame centrale de la video la plus longue, pour chaque modalité. Il s'agirait d'inspecter ~60 images contenant chacune 5 lignes (pour 5 participants), et 4 colonnes (pour les 4 modalités GoPros + Tobii), et de collecter les numéros de vignettes e.g. des GoPros à l'envers. Puisque cette partie à plutôt trait à l'étape de la collection des données plutôt qu'a l'analyse, je propose que vous vous occupiez de cette étape d'inspection. Au delà du fait que cela permette de se familiariser un peu avec les données, vous pourriez avoir de bonne idées sur d'éventuels statistiques à retenir pour la présentation des données ou le nettoyage du dataset (qu'il serait avantageux de collecter lors de cette étape d'inspection visuelle).  Par exemple, pour la GoPro2 du Gateau, bien que la majorité des videos soient orientées vers l'évier/le plan de travail, un certain nombre sont en buté en direction du fond de la piece. Pensez-vous qu'il faudrait qu'on rapporte la proposition de GoPro2 affectés ? Pour l'analyse, je pense que celle-ci seront plutôt inexploitable dans un premier temps et donc il faudrait aussi les répertorier.

% Dans un premier temps, je vous propose de regarder l'échantillon de ces vignettes dans le dossier 'snapshots' sur Nextcloud (avec les slides de la presentations)*. Est-ce que vous pensez à d'autres informations importantes à collecter lors de cette inspection visuelle ?

%  Encoding

% 1. 'is_upside_down': bool =  Indicates whether videos are upside-down. 
% 2. 'is_swapped_modality' : bool = Indicates whether the videos are in the right folder. 
% 3. 'true_modality' : str = Name of the true modality if relevant. 
% 4. 'is_sink': bool = Unique to cuisine-GoPro3, describes wether the GoPro is in the sink extreme buttress. 
% 5. 'is_old_setup': For the cooking task, the old administrations had a chair in the angle of the room (instead of the lego table), and no hand sanitizers close to the sink, visible on GoPro2.
% 6. 

%  Procedure

% When opening a novel tile (with supposedly 25 thumbnails), first look at annotating:

% Go from very specific and fast to the most rare issues

% 1. the upside-down cases
% 2. the fish-eyed video This after some time help panticipating the other categories
% 3. Whether the Tobii corresponds to the correct task (cooking or lego)
% 4. The old setups 
% 5. the swapped videos
% 6. general search of visual issues

%  Notes

% - The GoPro3 in wrong buttress with fish-eyes seem to capture close to regular GoPro2 (see e.g. 449 and 453.

%  Screenshot

% See page Notion to show results and categories 





\section{Unsupervised estimation of symbolic representation of egocentric vision during an ecological tasks}



    % \subsection{Evolution of assignment distance sequences over increasing vocabulary size}
    % \label{apd:chapter_5_distance_fitting}

    %     % The evolution of assignment distance sequences over increasing vocabulary size, reported in \Cref{apd:chapter_5_distance_fitting}, shows That overall fitting error decreases as a result of higher model capacity,



    %     \begin{figure}[!htbp]
    %         \centering
    %         \includegraphics[width = 0.9\linewidth]{chapters/chapter 5/figures/chapter_5_distances_across_K.png}    
    %         \caption{\textbf{Assignment errors of participants' symbolic representations across ego-vision prototypes vocabulary size $K$}. For increasing model capacity, each subfigures shows in each row the assignments distance of a participant performing the chocolate cake cooking task. We verify that the overall fitting error decreases as a result of higher model capacity, suggesting higher vocabulary size is required to properly estimate participant's ego-vision ethograms with controlled error bounds. Importantly, the heterogeneous fitting error trajectories, which appear relatively consistent across participants, yields two important insights for the final model estimate. First, that some steps of the task ought to receive higher proportion of the model's parameters to improve local disambiguation of the ego-vision prototypes. This observation calls for careful modeling of the temporal structure  the temporal distributions of temporal neighbors in the task will play a crucial role in disambiguating  the discriminative power of in disambiguating temporal neighbors. % Could we also hypothesize that the further in the task the more naturally disambiguation is "easy" ? I don't think so because the first steps are the most similar across participants, and so the most difficult to disambiguate.}
    %         }
    %         \label{fig:chapter_5_distances_across_K}
    % \end{figure}


    \subsection{Assignment distance condition in the definition of the residual spaces}
    \label{apd:chapter_5_distance_distribution_thresholding}

        The residual space at step $p\in [P]$ aim to contain the egocentric vision atoms whose semantic is not already captured by current cumulated set of centroids. These visual concepts exhibit within-cluster variability observed empirically in the annotation frames and encoded in the cluster assignment distance distributions.
        In practice, these visual concepts are determined unequivocally from inspecting many independent samples of the nearest region of centroids influence areas. This conceptual consistency for furthest independent samples only remain for a portion of centroids - suggesting a low amount of false positive assignments - but disappear for a significant proportion of other clusters. We aim to address this by integrating these mis-clustered samples to future stages of centroids estimations. 

        We perform an empirical study of different thresholding approaches, including:
        \begin{itemize}
            \item $d_{min}^k = Q_1(d_k)$ where $Q1$ denotes the first quartile of cluster $k$ assignment distance
            \item $d_{min}^k = Q_2(d_k)$ where $Q2$ denotes the median of cluster $k$ assignment distance
            \item $d_{min}^k = \text{mean}(d_k) + \text{std}(d_k)$  one standard deviation over the empirical average assignment distance
            \item $d_{min}^k = \text{mean}(d_k) + 2 * \text{std}(d_k)$  two standard deviation over the empirical average assignment distance
            \item $d_{min} \in [0.05, 1]$ varying distance threshold. 
        \end{itemize} 

        where $ d_k = \{ d(x, \mu_k) \mid x \ in V_k \}$ is the distribution of assignment distance for centroid $k$. \Cref{fig:chapter_5_residual_space_threshold} shows the evolution of residual space samples percentage for different thresholding technics.

        \begin{figure}[!htbp]
            \centering
            \includegraphics[width=0.9\linewidth]{chapters/chapter 5/figures/chapter_5_resildual_space_threshold.png}    
            \caption{\textbf{Percentage of samples integrated to the residual space across thresholding techniques.} After each iteration of the procedure, samples assigned to a centroid with an assignment distance above a threshold $d_{min}^k$ are integrated to the residual space. Different thresholding technics are investigated, including fixed distance thresholds ($1^{st}$ quartile, median distance,  , and adaptive thresholds based on empirical cluster assignment distance distributions.}
            \label{fig:chapter_5_residual_space_threshold}
        \end{figure}

 \subsection{Post-hoc interpretability of video encoder latent representations}
    \label{apd:chapter_5_visual_probing}


    The identification of each cluster semantics from these subsets of participants requires an efficient and reliable method to inspect the corresponding video frames. The embeddings of VideoMAE-v2 represent a spatio-temporal content of about $640ms$ of video (16 frames), incorporating both appearance and motion information. 

    In theory, for $30$ participants, validating \textbf{H1} and \textbf{H2} would require the inspection of $2 \times 30 \times 16 = 960$ frames per cluster, which is time-consuming and impractical for large vocabularies. However, given typical head-pose dynamics and the short duration of each input sequence (16 frames, $\approx 640$\,ms), we hypothesize that the central frame suffices in most cases (\textbf{H3}).
    
    Without motion, the representation can be fully described with the appearance of a single frame (e.g., the middle one). The motion component introduces additional degrees of freedom for these representations, which may complicate recovering their semantics. In our controlled setting, the primary sources of motion in video frames are:
    \begin{itemize} 

        \item \textbf{Camera motion} resulting from the participant’s head movements, which induces global motion blur characterized by a dominant blur direction.

        \item \textbf{Object motion} caused by the participant’s actions (e.g., grasping an object or washing hands), leading to more complex, localized blur patterns.
    \end{itemize}
        % \footnotetext[1]{Video foundation model is trained to capture both appearance and motions within the representations Empirical post-hoc studies of video foundation models representation investigate these aspects in [find ref] What about for egocentric vision VFM ? }
    We evaluate here the hypothesis proposed in \Cref{sec:chapter_5_annotation}, that the central of the 16 frames composing the encoder input sequence can be used as a relevant visual probe to recover an embedding semantics. We observed the sequence of frames within the encoder temporal horizon for $200$ randomly sampled embeddings, and report illustrative samples of the different scenario we found in \Cref{fig:chapter_5_visual_probing}. 

    \begin{figure}[!htbp]
        \centering
        \includegraphics[width = 0.9\linewidth]{chapters/chapter 5/figures/chapter_5_visual_probing.png}    
        \caption{\textbf{Illustrative examples of eight equally spaced frames within the temporal horizon (16 frames; $\approx 640ms$) of the video encoder.}}
        \label{fig:chapter_5_visual_probing}
    \end{figure}

    In most cases, we observed that the visual appearances remained consistent (e.g., rows a-e in \Cref{fig:chapter_5_visual_probing}), suggesting that the inspection of the central frame is a relevant visual probe and is sufficient to recover clusters' semantics. However, head and hands motions can create situations where the central frame alone cannot convey the full semantics associated with the performed action (e.g., rows (f-i) in \Cref{fig:chapter_5_visual_probing}), which limits the validity of \textbf{H3}.  of our approach. Nevertheless, we observed that for motion-rich clusters that emerge during clustering, inspecting multiple independent central frames
    can still reveal consistent patterns of appearance and motion blur, making it possible to infer the action semantics and alleviating this limitation (see \Cref{fig:chapter_5_ambiguous_examples}-g and \Cref{fig:chapter_5_definitive_examples}-f).

    % In fact when we say inspection it encompass all what we mean by recovering the semantic meaning, etc. 
    % We now discuss the extent to which latent representations of $w=16$ frames can yield an  
    % The assessement of semantic patterns associated with these participants samples depends on our ability to inspect sample embeddings efficiently, as well as the type of information they encode. 

    \subsection{Model selection problem of estimating an optimal number of clusters}
    \label{apd:chapter_5_model_selection_G}
        Determining an optimal number of clusters in unsupervised setting is a classic model selection problem with a broad literature \cite{stop_using_elbow, HENNIG201553, Sugar2003}, including (but not limited to) unsupervised validation indices such as Silhouette \cite{ROUSSEEUW198753}, Calinski-Harabasz \cite{Calinski01011974}, and Davies-Bouldin \cite{Davies1979}; stability-based or resampling approaches (e.g., prediction strength \cite{Tibshirani2005}, consensus clustering \cite{10.1162/153244303321897735, Monti2003}, cluster stability \cite{BenDor1999}); information-theoretic and likelihood-based criteria (e.g., Bayesian information criterion (BIC; \cite{Schwarz1978}) or Akaike information criterion (AIC; \cite{Akaike1998}) for model-based clustering \cite{Fraley2002}); and reference-distribution methods such as the Gap Statistic \cite{Tibshirani2001} (see additional discussions in \cite{HENNIG201553}). As a standard method shown to perform well in practice, we used the silhouette score in our experiments. 

        \paragraph{Validity and limitations of the silhouette score.\\} 
        
        Despite its popularity, the silhouette score has recognized limitations. It is (i) sensitive in high-dimensional spaces, where distance concentration reduces contrast between clusters \cite{Aggarwal2001}, and (ii) biased towards spherical and non-overlapping geometries, potentially misjudging elongated or chain-like clusters where intra-cluster dissimilarity $a(k)$ grows while inter-cluster dissimilarity $b(k)$ remains stable. In our setting, these concerns are mitigated for two reasons. First, the computation operates on cluster's representative rather than all samples, thereby reducing sensitivity to cluster shape and noise, and  enhancing data quality during silhouette score computation compared to traditional settings \cite{Batool2021, Hennig2015, deAPMarques2022}. This differs from previous uses of silhouette maximization \cite{Batool2021, Hennig2015, deAPMarques2022}, where the metric evaluates the quality of empirical cluster assignments. Here, it instead assess the quality of the aggregation of single centroids into connected components to mitigate modest levels of redundancy, thereby enhancing data efficiency during silhouette score computation. 
        Second, centroids are already estimated by minimizing the $k$-means objective, which ensures compactness by reducing within-cluster variance. The silhouette score complements this criterion by additionally normalizing for separation, providing a metric that balances both cohesion and isolation of clusters \cite{Batool2021, Shutaywi2021}.



    \subsection{Flat clusters computations from a dendrogram of accumulated centroids}
    \label{apd:chapter_5_flattening}




    We detail here the estimation of $G$ flat clusters from the dendrograms representing the structure of our accumulated centroids (see \Cref{sec:chapter_5_prototyping}).
    The goal is to enforce a fixed number of clusters (\(G\)), to obtain a $G$-flat partition obtained by cutting the dendrogram at the level that produces exactly \(G\) disjoint clusters.    
     
    Let \(\mathcal{T}\) be the dendrogram produced by HAC using the multi-view dissimilarity \(D_M^{[K]}\). Each internal merge \(m\) occurs at a \emph{height} (linkage value) \(h(m)\), yielding a strictly increasing sequence of merge heights \(0 < h_1 < \dots < h_{K-1}\). For any threshold \(t \in [h_g, h_{g+1})\) (with the conventions \(h_0:=0\), \(h_K:=+\infty\)), removing all edges with height \(> t\) partitions the leaves into exactly \(G=K-g\) connected components, which induces a flat clustering (partition) \(\nu^{[G]}(t) \in \mathcal{P}(X)\). Equivalently, a clustering with $G$ clusters can be obtained by choosing the unique threshold \(t_G \in [h_{K-G},\, h_{K-G+1})\) that yields \(|\nu^{[G]}(t_G)|=G\).
    We denote the corresponding assignment by \(\Phi_{t_G} : [K]\to [G]\), with prototypes \(\nu_g(t_G) = \{\mu_k \mid k\in \Phi_{t_G}^{-1}(g)\}\).

    % \paragraph{Inconsistency criterion.\\}

    %     The second is to rely on the so-called \emph{inconsistency criterion}, which adaptively determines cluster boundaries by comparing the linkage distance at a given merge with the average and variance of linkage distances over the preceding \(d\) levels in the subtree.            


    %     Formally, the inconsistency coefficient at node \(j\) is defined as
    %     \[
    %     I_j(d) \;=\; \frac{\ell_j - \mu_j(d)}{\sigma_j(d)},
    %     \]
    %     where \(\ell_j\) is the linkage distance at merge \(j\), and \(\mu_j(d), \sigma_j(d)\) are respectively the mean and standard deviation of linkage distances among the last \(d\) merges below \(j\). 
    %     A cut at threshold \(T\) then splits the tree wherever \(I_j(d) > T\). 
    %     This introduces two hyperparameters: the depth \(d\), controlling whether inconsistency is evaluated locally or more globally, and the threshold \(T\), determining the granularity of the cut. 

    % \paragraph{Experimental setup.\\}

    %     In practice, we explored a grid of thresholds (\(T \in [0.5,3]\)) and depths (\(d \in \{2,5,10\}\)), and compared the resulting partitions in terms of (i) number and size distribution of clusters, (ii) internal validation indices, and (iii) semantic interpretability of prototypes. 
    %     We observed that the inconsistency-based cuts were often more \emph{adaptive}, in the sense that they respected local heterogeneity in the dendrogram and yielded different partitions than fixed-\(G\) cuts with the same cardinality. 
    %     This adaptivity is expected in our setting, where semantic and visual similarity are not strictly transitive across subtrees (e.g., chopping vegetables is internally consistent but not interchangeable with pouring liquids, even though both belong to the broader class of preparation activities). 



    \subsection{Comparison of clustering algorithms performances}
    \label{apd:chapter_5_clustering_technics}
    % We use cosine over euclidean metric and Faiss K-Means over Scikit-learn implementation,
    We compare the performances of two K-Means implementations: Faiss \cite{douze2024faiss} and Scikit-learn \cite{sklearn}, using both cosine and euclidean distance metrics. We report in \Cref{tab:quality_counts} the number of valid and ambiguous centroids obtained with each combination of distance metric and implementation, for a vocabulary size of $C=100$. 

    \begin{table}[!htbp]
        \centering
        \begin{tabular}{lcc|cc}
            \hline
            \multirow{2}{*}{\textbf{Category}} & 
            \multicolumn{2}{c|}{\textbf{Cosine}} & 
            \multicolumn{2}{c}{\textbf{Euclidean}} \\
            & \textbf{Faiss} & \textbf{Scikit-Learn} & \textbf{Faiss} & \textbf{Scikit-Learn} \\
            \hline
            Valid & 75 & 69 & 68 & 54 \\
            Ambiguous & 25 & 31 & 32 & 46 \\
            \hline
        \end{tabular}
        \caption{Number of valid and ambiguous centroids across distance metric and implementation ($C=100$, $p=1$).}
        \label{tab:quality_counts}
    \end{table}

    We observe that Faiss yields a higher number of valid centroids, regardless of the distance metric used. Additionally, we find that the cosine distance metric outperforms the euclidean distance in both implementations. These results support our choice of using Faiss with cosine distance for our clustering approach. We leave as future work the exploration of other clustering algorithms, such as hierarchical clustering or DBSCAN, to further improve the clusters' quality and ultimately the convergence rate of our iterative exploration algorithm.


    \subsection{Related work on temporal action segmentation (TAS)}
    \label{apd:chapter_5_utas_related_work}

        % TAS commonly relies on two core components: (i) frame or video snippet level representations, which extracts informative features capturing spatial appearance and short motion information, and (ii) temporal and sequential modeling, which improves incorporates temporal dependencies and sequential context for improved TAS accuracy \cite{TAS_survey} 
    
    %We present in \Cref{apd:chapter_5_utas_related_work} related work addressing the TAS.

    Supervised approaches to temporal action segmentation (TAS) have achieved significant progress in recent years\footnotemark[1] due to the introduction of large-scale benchmarks (e.g. Ego4D \cite{ego4D}, EPIC-KITCHENS \cite{epic}); powerful video backbones for extracting visual motion features from RGB frames and optical flow using convolutions (Inflated 3D ConvNET (I3D) \cite{i3d}) or transformers-based architecture (e.g. videoMAE-V2 \cite{videomae-model}); and increasingly sophisticated temporal reasoning mechanisms. A recent survey \cite{TAS_survey} categorized them broadly as: (i) frame-level approaches, which learn temporally-aware feature representations by extending the receptive field of neural architectures (e.g., recurrent neural networks \cite{9156462}, temporal convolutional networks \cite{8099596, li2020ms}, and Transformers \cite{chinayi_ASformer, 10.5555/3295222.3295349}); and (ii) segment-level approaches, which explicitly model the sequential dynamics of actions, for instance with Hidden Markov Models and Viterbi decoding \cite{8578869, kukleva2019unsupervised, 6909500, 1054010}, Mallows Models modeling actions distributions over orderings or permutations \cite{sener2018unsupervised, MALLOWS1957}, and hierarchical Bayesian model \cite{goel2018learning} focusing on the procedural temporal structure of human activity. 
    \footnotetext[1]{See recent surveys for comprehensive reviews \cite{TAS_survey, 9839464}.}

    %Joint Visual-Temporal Embedding for Unsupervised Learning of Actions in Untrimmed Sequences:For instance, in videos showing how to make pancakes the process of cracking eggs should not only look visually similar across different videos but would generally occur before other tasks such as mixing the eggs into the batter or pouring the batter onto the griddle. So, it can be assumed that actions, at least in task-oriented videos, not only share certain visual features but also occur in a similar temporal space.

    However, these approaches often require unscalable and costly collection of high-quality annotations, supporting the recent development of promising weakly supervised and unsupervised solutions. Two main strategies are typically employed: clustering the feature space across videos to identify actions, mainly via joint visual-temporal self-supervised representation learning trained on one or a combination of visual, temporal, or sequential (e.g. ordering) predictive subtasks \cite{https://doi.org/10.48550/arxiv.2001.11122, 9157452, https://doi.org/10.48550/arxiv.1904.04189, Bojanowski2014} and/or using temporal positional encoding \cite{deAPMarques2022, retrocausal}; or primarily focusing on detecting action boundaries within single videos \cite{9879786, twfinch, 990935}. Additional information discussing the related work can be found in \Cref{apd:chapter_5_utas_related_work}. Altogether, unsupervised approach have demonstrated remarkable performances sometimes outperforming supervised approaches, largely by exploiting the inherently strong spatio-temporal and sequential regularities in human activity videos. 

    Given the data-driven prototyping approach in \Cref{sec:chapter_5_prototyping} focuses on modeling interpretable action classes, we propose a complementary strategy: detecting action boundaries by exploiting temporal patterns in the latent similarity sequences directly, promoting the transferability and interpretability of our methodology.                 

    As discussed in the preliminary analysis in \Cref{sec:chapter_4_preliminary_analysis_egocentric_latent_sequences}, the temporally ordered sequences of video snippet embeddings exhibit key temporal patterns. These include the \textit{diagonal blocks} with consistently elevated cross-samples similarity - suggesting an action segment - and the short \textit{cross-shaped artifacts} - suggesting transitions between actions. These patterns are explicitly leveraged in the state-of-the-art unsupervised approaches: TW-FINCH demonstrates that temporal progression can suffice to discover actions via hierarchical clustering of temporally weighted 1-nearest-neighbor graph \cite{twfinch}.  In contrast, Action Boundary Detection (ABD) relies on the contrast between intra-segment similarity and inter-segment dissimilarity, applying non-maximum suppression on smoothed similarity profiles to detect change-points \cite{9879786}. While effective, both approaches present limitations: (i) they assume the number of action classes is known to refine their change-points estimate (with ABD proposing an online version using a constant threshold on the non-maximum decision to withdraw pseudo boundaries), and (ii) overlook specificities of egocentric video exhibiting larger amount of pseudo-boundaries due to head motion \cite{sam_icassp}, hindering their efficiency for egocentric content.
    %Triplet loss \cite{10209008}
    %ABD \cite{9879786}

    We propose instead to exploit these temporal patterns via the optimization of a multiple change-point objective using reproducing kernels (KCP) \cite{cpt}, which presents the significant benefits of (i) being theoretically grounded, with established consistency \cite{garreau2018consistent} and finite-sample optimality results under reasonable assumptions, taking the form of an oracle inequality in terms of quadratic risk \cite{CELISSE2018200, cpt}; (ii) nonparametric models (i.e., inference-ready and independent of stored parameters or training data), in particular the procedure do not require to know the number of change-points in advance; and (iii) presenting efficient optimization algorithms (e.g. PELT \cite{pelt}) with fast implementations (\texttt{ruptures} package \cite{ruptures}). These were driven by key findings from the literature, which is described below along with our methodology.





    % Most work on temporal action segmentation (TAS) is framed as a supervised, deep-learning problem, typically with strong frame encoders and temporal models (e.g., Temporal Convolutional Network (TCN) and Transformers) trained on dense labels \cite{TAS_survey}. Yet unsupervised variants are gaining traction as video collection expands, especially long, untrimmed egocentric streams, where exhaustive annotation is costly. Large benchmarks like Ego4D \cite{ego4D} and EPIC-KITCHENS \cite{epic} were proposed to advance the field. %TAS surveys typically organize methods by supervision level (fully/weakly/unsupervised), representation learning, temporal modeling, and post-processing;
    % Surveys focusing on egocentric content emphasize hand–object interaction, gaze and attention, object-state change, and multi-modal cues that motivate label-efficient segmentation pipelines \cite{NUNEZMARCOS2022175}. Unsupervised TAS typically follows one of a few directions: (i) optimal-transport decoding that imposes temporal consistency by fusing a Gromov–Wasserstein prior with unbalanced OT to produce coherent segments and then self-train on the pseudo-labels \cite{Xu_2024_CVPR}; (ii) boundary-first changepoint detection via frame-similarity with non-maximum suppression (NMS) and light segment clustering \cite{9879786}; (iii) self-supervised learning that discovers a discriminative space and applies generic clustering—e.g., a shallow triplet-loss framework; and (iv) “exotic” but effective object-centric boundary detectors that fuse local and global features—e.g., OTAS. Together, these lines reduce annotation dependence while remaining competitive on long, untrimmed instructional/egocentric videos.    

    % Unsupervised action segmentation aims to automatically detect temporal actions boundaries in videos and classify them into action categories (including background segments that do not contain relevant high-level actions), which is needed for various tasks related to the understanding of human activity in videos \cite{9839464}.

    % Supervised approaches have received a lot of attention to address temporal segmentation of videos. For instance, transformers-based models \cite{Zhang_Wu_Li_2022} and two-stream temporal convolutional networks \cite{mstcn} have been used to model long-range time dependency across the video, potentially post-processing the output segmentation using graph-based models \cite{9156462}.
    % Weakly-supervised approaches have been proposed to mitigate the cost of annotations and their subjectivity and use either the ordered sequence of actions \cite{https://doi.org/10.48550/arxiv.1805.06875}, sparse annotations e.g. labelling only an arbitrary frame within each action segment  \cite{https://doi.org/10.48550/arxiv.2103.06669}, or human-generated information such as speech or captions \cite{https://doi.org/10.48550/arxiv.1506.09215}.

    % However, these approaches follow a paradigm that relies heavily on training data, which limits their practical value and requires high-quality image-level annotations at the cost of significant, time-consuming and unscalable human effort. 

    % To circumvent this problem,  unsupervised approaches have recently emerged. Some work uses alternate self-supervised learning of visual and temporal appearance representations of the frames of multiple videos of the same activity before clustering them to obtain the action segmentation \cite{https://doi.org/10.48550/arxiv.2001.11122}. These two steps can be merged to learn an embedding space while jointly clustering the frames representation using a temporal optimal transport module that preserves detected actions ordering \cite{retrocausal}. In the same spirit, actions in multiple videos can be co-localized by iteratively optimizing a clustering of the video frame representation and the localization of action boundaries \cite{Gong_2020_CVPR}. However, these works make the restrictive assumption that several videos of the same activity are available. A state-of-the-art approach, TW-FINCH, overcomes this limitation by building a time-weighted  nearest-neighbor graph from the frame embedding, and discovers action segments by applying a recursive hierarchical clustering algorithm. Connected components of the graph form clusters with similar appearances \cite{twfinch}. However, it is assumed that the camera is fixed, which might limit their performances on egocentric content. 


    % \cite{10483653}: permutation aware ...


    % Another recent weaky supervised work collect only a single timestamps per action and perform random walks on graphs to perform the temporal segmentation \cite{Hirsch2024}.


    % The selection of a suitable penalization is a classic but challenging model selection problem \cite{Lavielle2005,Brodsky1993}, with methods including popular criteria (AIC \cite{Akaike1998}, Mallow's Cp \cite{Mallows1973}, BIC \cite{Schwarz1978}), $k$-fold cross validation \cite{Golub1979}, and stability analysis \cite{Liu2010, Meinshausen2010}. A complete discussion and additional references can be found in \cite{cpt}. 

    
    % is unsupervised and provides additional theoretical justification for the penalization term, originally found in \cite{cpt} (see in \cite{CELISSE2018200}). Their method is presented in \Cref{alg:KCP}.    

    


    %  \Cref{fig:chapter_5_slope_heuristic}

    % \begin{figure}[!htbp]
    %     \centering
    %     \includegraphics[width = 0.9\linewidth]{chapters/chapter 3/figures/chapter_5_slope_heuristic.png}
    %     \caption{\textbf{Estimation of the kernel change points detection penalty parameters using the slope heuristics} Following \cite{slop_heuristics_survey}, we perform a linear regression of $ \big(\hat{\mathcal{R}}_n(\hat{\tau^{(M)}}) \big)_{m \in [0.6M_{\text{max}}, M_{\text{max}}]     }$ }}
    %     \label{fig:chapter_5_slope_heuristic}
    % \end{figure}




    \subsection{Definitions of clustering metrics used to compare alignment of symbolic sequences}
    \label{apd:chapter_5_clustering_metrics_definitions}

    To compare the alignment of the symbolic representations obtained with different vocabularies, we use two standard clustering metrics: the Adjusted Rand Index (ARI) \cite{Hubert1985} and the Normalized Mutual Information (NMI) \cite{Lancichinetti2009}. These metrics quantify the similarity between two clusterings or sets of prototypes, providing insights into how well different vocabularies capture the underlying structure of the data. We provide below their formal definitions.

    \begin{dfn}
        \textbf{Adjusted Rand Index (\emph{ARI}) \cite{Hubert1985}.}
        Consider two clusterings, \(S_1\) and \(S_2\), of a dataset with \(n\) elements. To compare these clusterings, we count the number of samples that are:
        \begin{itemize}
            \item \text{TP} (\emph{True Positive}): in the same cluster in both \(S_1\) and \(S_2\),
            \item \text{FP} (\emph{False Positive}): in the same cluster in \(S_1\) but in different clusters in \(S_2\),
            \item \text{FN} (\emph{False Negative}): in different clusters in \(S_1\) but in the same cluster in \(S_2\),
            \item \text{TN} (\emph{True Negative}): in different clusters in both \(S_1\) and \(S_2\).
        \end{itemize}

        The \emph{Rand Index} (RI) measures the proportion of samples that agree between the two clusterings:
        \[
            \text{RI}(S_1, S_2) = \frac{\text{TP} + \text{TN}}{\text{TP} + \text{FP} + \text{FN} + \text{TN}}.
        \]

        However, the RI does not account for chance agreement. The \emph{Adjusted Rand Index} (ARI) corrects for this by normalizing the RI with respect to its expected value under random clusterings:
        \[
            \text{ARI}(S_1, S_2) = \frac{\text{RI}(S_1, S_2) - \mathbb{E}[\text{RI}(S_1, S_2)]}{1 - \mathbb{E}[\text{RI}(S_1, S_2)]},
        \]
        where \(\mathbb{E}[\text{RI}(S_1, S_2)]\) is the expected Rand Index for random cluster assignments.

        The ARI ranges from \(-1\) (indicating worse than random agreement) to \(1\) (perfect agreement), with \(0\) indicating agreement equivalent to random chance.
    \end{dfn}

    \begin{dfn}
        \textbf{Normalized Mutual Information (\emph{NMI}) \cite{Lancichinetti2009}.}
        The Normalized Mutual Information (NMI) quantifies the mutual dependence between two clusterings or sets of prototypes, normalized by their entropies. Given two discrete random variables \(S_1\) and \(S_2\) representing cluster assignments (or symbolic sequences), the NMI is defined as:
        \[
            \mathrm{NMI}(S_1, S_2) = \frac{2 \cdot I(S_1; S_2)}{H(S_1) + H(S_2)},
        \]
        where:
        \begin{itemize}
            \item \(I(S_1; S_2)\) is the \emph{mutual information} between \(S_1\) and \(S_2\), measuring the amount of information obtained about one variable through the other,
            \item \(H(S_1)\) and \(H(S_2)\) are the \emph{Shannon entropies} of \(S_1\) and \(S_2\), respectively, quantifying the uncertainty in each clustering.
        \end{itemize}
        The NMI ranges from \(0\) to \(1\), where:
        \begin{itemize}
            \item \(1\) indicates perfect correspondence between the two clusterings,
            \item \(0\) indicates no statistical relationship.
        \end{itemize}
    \end{dfn}

    
    \subsection{Additional information on prototypes}
    \label{apd:chapter_5_prototype_additional}

    \paragraph{Histograms of assignment distance}

        % 1) discuss the prototypes distances distributions 
        As illustrated in \Cref{fig:chapter_5_distance_distributions}, the assignment distance distributions vary significantly across centroids, with some prototypes exhibiting a high degree of semantic consistency (i.e. low assignment distances) while others exhibit a more diverse set of samples (i.e. high assignment distances).

        \begin{figure}[!htbp]
            \centering
            \includegraphics[width = 0.9\linewidth]{chapters/chapter 5/figures/chapter_5_distance_distributions.png}    
            \caption{\textbf{Randomly selected histograms of clusters assignment distance.} The assignment of empirical samples to their nearest prototypical region exhibit various patterns of assignment, as suggested by the prevalence, skewness, and average assignment distances.}
            \label{fig:chapter_5_distance_distributions}
        \end{figure}


        

        %CODE: TODO
        % \paragraph{Vocabulary comparison based on the unsupervised symbolic representation indicators.\\}
        % \label{sec:chapter_5_comparison_across_models}
        % % Introductory sentence ahta is more high level reminding the link between gram matrices diagonal blocks and the prototyping granularity

        % % Motifs length between raw and after segment level consensus 

        % To evaluate whether the learned prototypes achieve the desired representational granularity, we study the relationship between estimated action boundaries and prototype assignments. The segmentation step identifies change-points in the reproducing kernel Hilbert space as “mean-change” events (see \Cref{sec:chapter_5_unsupervised_temporal_segmentation}). We thus compute the proportion of \textit{pseudo change-points}, i.e., estimated change-points that occur between two consecutive identical prototype labels.

        % This indicator highlights potential mismatches between the temporal segmentation and the prototyping modules: frequent pseudo change-points suggest that prototypes cover larger temporal scales than the intended fine-grained granularity. However, this measure may be misguiding if (i) the prototype vocabulary lacks exhaustivity, such that truly distinct visual actions are collapsed into a single prototype (epistemic uncertainty), and (ii) when the underlying subsequences exhibit strong variability and cannot be mapped to piecewise-constant dynamics (e.g., highly changing visual experience; aleatoric uncertainty).

        % \paragraph{Analysis of false positive assignments}
        %     False assignment of samples to the prototypes may occur as a result of: 

        %     -  during action transition phases, 
        %     - the realistic lack of exhaustivity for fine-grained prototyping occurring at such short spatiotemporal scale, i.e. samples may be misrepresented
        %     -  or imperfect sample encoding,



\subsection{Additional examples of prototypes nearest- and furthest-neighbor visual probes}
\label{apd:chapter_5_visual_probes_additional}

    We provide in \Cref{apd:chapter_5_visual_probes_additional} additional examples of nearest- and furthest-neighbors visual probes for randomly selected prototypes from \texttt{Model-B}. These examples complement the illustration of \Cref{fig:chapter_5_visual_probes}, supporting the generalizability of the prototypes to unseen participants, as well as their semantic consistency within prototypical regions.

    \begin{figure}[!htbp]
        \centering
        \includegraphics[width=0.9\linewidth]{chapters/chapter 5/figures/chapter_5_visual_probes_second.png}    
        \caption{\textbf{Randomly selected visual probes of egocentric video prototypes.} We assess prototypes generalizability to testing participants via nearest-neighbors visual probes (left and middle columns). The semantic consistency within prototypical regions is assessed in testing participants by comparing nearest-neighbors and furthest-neighbors visual probes (middle and right columns).}
        \label{apd:chapter_5_visual_probes_second}
    \end{figure}


\subsection{Additional examples of valid and rejected prototypes}
\label{apd:chapter_5_prototype_visual_probe_additional}
    

    To complement the illustration of our vocabulary consolidation step (\Cref{sec:chapter_5_hierarchical_multiview_clustering}), we provide additional examples of accepted and rejected prototypes in \Cref{fig:prototype_visual_probe_additional_accepted} and \Cref{fig:prototype_visual_probe_additional_accepted}, respectively.

    \begin{figure}[!htbp]
        \centering
        \includegraphics[width=0.7\linewidth]{chapters/chapter 5/figures/prototype_visual_probe_additional_accepted.png}
        \caption{\textbf{Example of an accepted prototype based on nearest-neighbors visual probes.} Each item corresponds to the nearest-neighbors visual probe of a centroid (see \Cref{sec:chapter_5_annotation}). In this example, the prototype is accepted, as all constituent centroids consistently capture the same action of "pouring the dough in the pan".}
        \label{fig:prototype_visual_probe_additional_accepted}
    \end{figure}

    \begin{figure}[!htbp]
        \centering
        \includegraphics[width=0.7\linewidth]{chapters/chapter 5/figures/prototype_visual_probe_additional_rejected.png}    
        \caption{\textbf{Example of a rejected prototype based on nearest-neighbors visual probes.} Each item corresponds to the nearest-neighbors visual probe of a centroid (see \Cref{sec:chapter_5_annotation}). In this example, the prototype is rejected, as the constituent centroids do not consistently capture the same underlying action, mixing "opening the fridge", "setting up the oven", and "looking at the kitchen counter".}
        \label{fig:prototype_visual_probe_additional_rejected}
    \end{figure}

    As can be seen, rejected prototypes may mix distinct actions that share visual attributes. In contrast, the accepted prototypes exhibit reasonable semantic consistency, with all constituent centroids representing the same action. This highlights the importance of careful prototype evaluation to ensure meaningful and interpretable representations of egocentric vision behavior.

\subsection{Additional examples of prototypes matching across vocabularies}


    \begin{figure}[!htbp]
        \centering
        \includegraphics[width=0.9\linewidth]{chapters/chapter 5/figures/chapter_5_matches_second_best.png}
        \caption{\textbf{Example of top-2 prototype matches in \texttt{Model-B} of the "adding the butter" \texttt{Model-A} prototype, based on the number of overlapping embeddings.}} %We observe that a marginal amount of samples assigned to "reading the recipe" using \texttt{Model-A} were assigned to a second best match (right) that does not match with the anchor prototype semantics.}
        \label{fig:chapter_5_matches_second_best}
    \end{figure}

    \begin{figure}[!htbp]
        \centering
        \includegraphics[width=0.9\linewidth]{chapters/chapter 5/figures/chapter_5_best_match_recipe.png}
        \caption{\textbf{Example of top-2 prototype matches in \texttt{Model-B} of the "reading the recipe" \texttt{Model-A} prototype, based on the number of overlapping embeddings.}}
        \label{fig:chapter_5_best_match_recipe}
    \end{figure}


\section{Symbolic representation analysis and barycenter estimation for the assessment of dysexecutive syndromes}

    \subsection{Comparison of symbolic representations composition between control and patients groups}
    \label{apd:chapter_5_control_vs_patients}

        We provide in \Cref{fig:per_label_group_p_values} additional information on the comparison of symbolic representations composition between control and patients groups, complementing the analysis presented in \Cref{sec:chapter_5_symbolic_representation}. Specifically, we report the number of occurrences, average duration, and proportion of time for each labelled category, along with statistical significance of pairwise group differences within each category.

        \begin{figure}[h!]
            \centering
            \includegraphics[width=1.0\textwidth]{chapters/chapter 6/figures/per_label_group_p_values.png}
            \caption{\textbf{Comparison of number of occurrences (top), average duration (middle), and proportion of time (bottom) of each labelled category between control and patient diagnosis groups.} Bars indicate mean values per category and group, with participant counts displayed above or below each bar. Error bars represent standard deviation divided by the square root of the sample size (standard error of the mean). Pairwise group differences within each category were tested using the Mann-Whitney U test with Bonferroni-Hochberg correction, and significant differences are shown by brackets connecting the bars with annotated p-values. Statistical significance was set at $p<0.05$.}
            \label{fig:per_label_group_p_values}
        \end{figure}


    \subsection{Analysis of barycenters including the background symbol}
    \label{apd:chapter_6_barycenters_background}

        We provide in \Cref{fig:chapter_6_bg} the estimated barycenters for each diagnosis groups when including the background symbol when updating the barycenter at each iteration. While the overall patterns remain similar to those presented in \Cref{fig:chapter_6_non_bg}, we observe that the inclusion of the background symbol leads to more fragmented barycenters, with a higher prevalence of the background symbol interrupting sequences of action prototypes. 
            \begin{figure}[!htbp]
                \centering
                \includegraphics[width = 1.0\linewidth]{chapters/chapter 6/figures/chapter_6_bg.png}    
                \caption{\textbf{Estimated barycenters for each diagnosis group.} Each barycenter is computed using the DBA algorithm under the TW-TWE distance. Barycenters are estimated including the background symbol when updating the barycenter at each iteration.}
                \label{fig:chapter_6_bg}
            \end{figure}

    \subsection{Analysis of barycenters using the Euclidean distance}
    \label{apd:chapter_6_barycenters_median_twe}

        We present in \Cref{fig:chapter_6_barycenters_median_twe}-a the estimated barycenters for each diagnosis groups using the TWE distance, illustrating the importance of inducing a meaningful distance measure between symbols for computing barycenters. Additionally, we present in \Cref{fig:chapter_6_barycenters_median_twe}-b the estimated barycenters using symbols majority voting after interpolating symbolic sequences to the median length. This approach yields degenerated barycenters, as it does not account for elastic alignment between sequences, leading to a loss of temporal dynamics information.

        \begin{figure}[!htbp]
            \centering
            \includegraphics[width = 1.0\linewidth]{chapters/chapter 6/figures/chapter_6_barycenters_median_twe.png}    
            \caption{\textbf{Estimated barycenters for each diagnosis group using the TWE distance (a) or the median (b).} The TWE distance incorporates an arbitrary ordering of the symbols distance, yielding unfaithful average barycenters. Likewise, the median approach does not account for elastic alignment between sequences, leading to temporally inconsistent barycenters.}
            \label{fig:chapter_6_barycenters_median_twe}
        \end{figure}



% \subsection{Procedure stopping criteria}
% \label{apd:chapter_5_stopping_criteria}

% Accounting for all participants, this number varies from $N=223 (STD=74)$ for raw timestamps estimates to $N=202 (STD=65)$ when combining samples labels estimates to reduce oversegmentation. In the experiments, the procedures is stopped at round $P=8$ with 


% Stopping criteria may include stopping when the residual set at step $p$ is small enough, e.g. $\| R_p \| \leq 0.1N_{train}$. In our experiments, the proportion of residual samples in the first (respectively in the last round $P=8$) covered $66.8\%$ (respectively $58.2\%$) of the training set. 

% At each iteration, the human coding step allows to assess at least subjectively levels of novelty, redundancy, and overall clustering qualitative performances. We observed that while novel vision concepts emerge frequently during the first few iterations of the procedure, suggesting an efficient residual set sampling, this effect seemed to decrease after round $p\approx 6$. This suggest that 



% - the first quartile of distances distributions
% one and two  the avarage distance 
% forthresholding technics to sampled  to c  ths 


% When observing independent samples in the nearest and furthest regions of the centroids influence areas, the degree of semantic associations with the reference unequivocal  varies from highly correla

% between samples vary from highly conceptual consictencies of encoded actions vary from highly consistent to  

% Examination of the furthest assessment of fursthest    samples are assigned to the closest of these visual concepts, 
% Here we remind the definition of the 



% \subsection{Clustering evaluation metrics}
% \label{apd:chapter_5_clustering_metrics}


%     We have considered several standard metrics to evaluate the quality of the meta-syllabes clustring across vocabulary size.

%     \begin{enumerate}
%         %TODO
%         \item \textbf{Silhouette score}~\cite{ROUSSEEUW198753}: The silhouette score is a measure of how similar an object is to its own cluster compared to other clusters. It ranges from -1 to 1, with higher values indicating better-defined clusters.
%         \item \textbf{Davies-Bouldin index}~\cite{davies1979cluster}: The Davies-Bouldin index is a metric that evaluates clustering algorithms based on the ratio of within-cluster distances to between-cluster distances. Lower values indicate better clustering.
%         \item \textbf{Calinski-Harabasz index}~\cite{calinski1974dendrite}: The Calinski-Harabasz index measures the ratio of the sum of between-cluster dispersion to within-cluster dispersion. Higher values indicate better clustering.
%         \item \textbf{Adjusted Rand index}~\cite{hubert1985comparing}: The Adjusted Rand index is a measure of the similarity between two data clusterings. It ranges from -1 to 1, with higher values indicating better agreement between the two clusterings.
%         \item \textbf{Normalized mutual information}~\cite{strehl2002relationship}: Normalized mutual information is a measure of the amount of information shared between two clusterings. It ranges from 0 to 1, with higher values indicating better agreement between the two clusterings.
%         \item \textbf{Modularity}~\cite{newman2004finding}: Modularity is a measure of the strength of division of a network into modules (or clusters). Higher values indicate a stronger division.
%         \item \textbf{Fowlkes-Mallows index}~\cite{fowlkes1983method}: The Fowlkes-Mallows index is a measure of the similarity between two clusterings. It ranges from 0 to 1, with higher values indicating better agreement between the two clusterings.
%         \item \textbf{Purity}~\cite{zhao2005evaluation}: Purity is a measure of the extent to which clusters contain only members of a single class. Higher values indicate better clustering.

%     \end{enumerate}


% CODE: TODO
% \subsection{Label noise injection near actions boundaries}
% \label{apd:chapter_5_boundaries_experiment}

% As mentioned in \Cref{sec:chapter_5_symbolic_representation}, the ambiguity of action transitions phase is often identified as a source of noise in contrast with the more stable regions of action segments, exhibiting more consistent similarity. Consistently with observation \Cref{obs:observation_similarity}, the degree to which empirical samples ressemnble prototypes is quantified by the assignment distance - here the cosine distance. To test whether near-boundaries samples differ in their average assignment distance compared to central onces, we tested one-sample XXX.

% quantify the assignment distance between prototypes and empirical samples is used to quantify the   source of noise  the 

% While it is expected from $k$-means to promote centroids within long and stable similarity consistent segments, this suggests that the boundaries constitute structurally less reliable regions for prototype assignment. 

% This is related to our prototyping procedure that promote the emergence of reflecting similarity-consisttent regions Likewise, it also suggests that centroids emerging from $k$-means objective tend to promote the prototyping of regions with stable similarity  consistent and  the iterative procedure arise or that prototypes emerge .


% \subsection{Matching of egocentric vision prototypes across models}
% \label{apd:chapter_5_assignment_model}



%     DFrom \cite{deAPMarques2022}
%     The problem of finding the mapping between the Ground-Truth
%     (GT) and the produced clusters can be regarded as an assignment
%     problem, as represented in Figure 2, in which the task is to find
%     a one-to-one mapping that maximizes the weights of a bipartite
%     graph where each edge is weighted by the number of matching
%     frames each GT action has within each cluster. In this particular
%     case, it is easy to manually produce the mapping that maximizes the
%     weights. To generate such a mapping, we follow the approach used
%     in [30, 45, 47] that is based on the Hungarian Method [28, 29, 38],
%     which tries to find a perfect matching: the one that matches all
%     vertices of the graph while maximizing its weights.
%     `'